\documentclass[11pt]{article}
\usepackage[margin=1in]{geometry}
\usepackage{amsmath,amssymb}
\usepackage{booktabs}
\usepackage{array}
\usepackage{enumitem}
\usepackage{hyperref}
\usepackage{xcolor}

\hypersetup{
    colorlinks=true,
    linkcolor=blue!60!black,
    urlcolor=blue!60!black
}

\title{\textbf{Graph Isomorphism by Continuous Relaxation and Hyperplane Rounding}}
\author{}
\date{Presentation draft, 2026-05-01}

\begin{document}
\maketitle

\begin{abstract}
This report documents the current codebase-grounded Phase 1 graph
isomorphism pipeline.
The method relaxes the discrete permutation search to a quadratic program over
the Birkhoff polytope, adds a degree-profile cost to the adjacency
commutator objective, rounds the relaxed matrix with a hyperplane method, and
accepts an isomorphism only when the rounded permutation passes the exact
integer check $AP = PB$.

The central result of Phase 1 is not that the method beats mature exact
isomorphism tools.
It does not.
The useful result is narrower: the relaxation gives an interpretable soft
correspondence matrix, the degree-profile term improves separation over an
adjacency-only objective, and the current bottleneck is now clearly located in
rounding quality and solver scaling.
\end{abstract}

\tableofcontents
\newpage

\section{Problem Statement}

Given two simple undirected graphs $G_A$ and $G_B$ with $n$ vertices, let
$A,B \in \{0,1\}^{n \times n}$ be their adjacency matrices in deterministic
node order.
The graphs are isomorphic if there exists a permutation matrix
$P \in \{0,1\}^{n \times n}$ such that
\begin{equation}
    AP = PB.
\end{equation}

Equivalently,
\begin{equation}
    P^\top A P = B.
\end{equation}

The direct search space contains $n!$ permutations.
The implemented approach replaces that discrete search with a continuous
relaxation, then tries to recover a discrete certificate.

\section{Relaxation}

A matrix $X \in \mathbb{R}^{n \times n}$ is doubly stochastic if
\begin{equation}
    X_{ij} \ge 0,\qquad
    \sum_j X_{ij}=1,\qquad
    \sum_i X_{ij}=1.
\end{equation}

The set of all doubly stochastic matrices is the Birkhoff polytope
$\mathcal{B}_n$.
Its vertices are exactly the permutation matrices.
This makes $\mathcal{B}_n$ the natural continuous relaxation of the discrete
permutation set.

In the solver, $X$ is interpreted as a soft node-correspondence matrix.
Large $X_{ij}$ means vertex $i$ in $G_A$ is being softly matched to vertex $j$
in $G_B$.

\section{Implemented Objective}

The current solver in \texttt{isomorphism.py} optimizes exactly two terms:
\begin{equation}
    \min_{X \in \mathcal{B}_n}
    \quad
    w_C \langle C, X \rangle
    +
    w_Q \|AX - XB\|_F^2.
    \label{eq:implemented-objective}
\end{equation}

The default objective weights are
\begin{equation}
    w_C = 1.0,\qquad w_Q = 1.0.
\end{equation}

No higher-order walk term, spectral alignment term, entropy regularizer, or WL
feature cost is currently active in the source-of-truth solver.
Those are future extensions, not present implementation.

\subsection{Adjacency Commutator}

The term
\begin{equation}
    Q(X)=\|AX-XB\|_F^2
\end{equation}
is the continuous version of the exact condition $AP=PB$.
If $X=P$ is a correct isomorphism, then $Q(P)=0$.

This is the cleanest baseline objective because it is directly tied to the
definition of graph isomorphism.

\subsection{Degree-Profile Cost}

The solver also builds a node-pair cost matrix from degree information.
Let
\begin{equation}
    d^A_i = \sum_k A_{ik}, \qquad d^B_j = \sum_k B_{jk}
\end{equation}
be vertex degrees, and let
\begin{equation}
    s^A_i = \sum_k A_{ik} d^A_k,\qquad
    s^B_j = \sum_k B_{jk} d^B_k
\end{equation}
be one-hop neighbor-degree summaries.

The raw cost is
\begin{equation}
    C_{ij}^{raw}
    =
    \alpha |d^A_i-d^B_j|
    +
    \beta |s^A_i-s^B_j|,
\end{equation}
with defaults
\begin{equation}
    \alpha=1.0,\qquad \beta=0.35.
\end{equation}

The code normalizes this cost matrix by its maximum absolute value when that
maximum is nonzero.
This matters because $\langle C,X\rangle$ and $\|AX-XB\|_F^2$ scale
differently with $n$.
Without normalization, one term can dominate for reasons unrelated to
isomorphism quality.

\subsection{Convexity}

The current objective is a convex quadratic program.
The adjacency term is the squared Frobenius norm of a linear map in $X$.
The degree-profile term is linear.
The feasible region $\mathcal{B}_n$ is convex.

This is important because the current solver does not require Gurobi's
\texttt{NonConvex=2} mode.
That would become relevant only after adding genuinely non-convex terms.

\section{Why Not Use Only $\|AX-XB\|_F^2$?}

The adjacency-only objective is mathematically clean, but practically weak.
In relaxed space, many fractional assignments can produce small adjacency
residuals without committing to a sharp vertex correspondence.

The degree-profile term supplies low-cost local information.
It discourages matching vertices whose degree and neighbor-degree summaries are
incompatible.
It does not solve graph isomorphism by itself, but it improves the shape of the
optimization landscape.

The ablation question for Phase 1 is therefore:
\begin{quote}
Does adding $\langle C,X\rangle$ improve separation between isomorphic and
non-isomorphic cases without increasing runtime materially?
\end{quote}

The current experiments support that direction.
For small reruns, isomorphic pairs still have objective values close to zero,
while verified non-isomorphic random pairs receive much larger objective values
when the degree-profile cost is enabled.

\section{Isomorphism Index}

After optimization, the solver reports the objective value $Z^*$ and a bounded
similarity score
\begin{equation}
    I = \exp(-\lambda Z^*),
\end{equation}
with default $\lambda=0.1$.

This index is useful for plots and comparisons.
It is not a proof.

The correct interpretation is:
\begin{itemize}[leftmargin=*]
    \item $I \approx 1$ means the relaxed objective found a very low-cost soft
    correspondence.
    \item $I \ll 1$ means the relaxed objective found substantial mismatch.
    \item Neither statement is an exact isomorphism decision.
\end{itemize}

\section{Rounding}

The relaxed solution $X^*$ is generally fractional.
To obtain a candidate isomorphism, the current solver calls
\texttt{hyperplane\_rounding.py}.

\subsection{Hyperplane Rounding Path}

The current rounding path is:
\begin{enumerate}[leftmargin=*]
    \item Compute an SVD embedding of rows and columns of $X^*$.
    \item Sample random hyperplanes in the embedding space.
    \item Convert projections into sign signatures.
    \item Build a Hamming-distance cost matrix between row and column
    signatures.
    \item Solve a linear assignment problem to obtain a permutation candidate.
    \item Repeat for multiple trials and keep the permutation with the smallest
    integer residual.
\end{enumerate}

The implementation runs 200 trials by default.
It stops early if any trial gives zero residual.

\subsection{Relation to Max-Cut}

The inspiration comes from Goemans-Williamson rounding for Max-Cut:
a continuous geometric solution is converted to a discrete object by random
hyperplanes.

Here the role is different.
We are not cutting one vertex set into two sides.
We are using random hyperplane signatures to convert the two SVD embeddings
from $X^*$ into a discrete matching problem.

\subsection{Hungarian Baseline}

Hungarian rounding remains an important comparison baseline.
Given $X^*$, it directly finds a high-weight permutation under the assignment
cost $1-X^*$.

In current small tests, hyperplane rounding has not yet shown a clear
certificate advantage over the Hungarian baseline.
That is not a failure of the report; it is a useful experimental result.

\section{Exact Verification}

The final decision is made only after rounding.
Given a candidate permutation $P^*$, the code checks
\begin{equation}
    AP^* = P^*B
\end{equation}
using integer-valued matrices.

If the equality holds, $P^*$ is a certificate of isomorphism.
If it does not hold, the result is not certified.

This is the key distinction:
\begin{equation}
    \text{failed certificate} \ne \text{proof of non-isomorphism}.
\end{equation}

The current method can prove isomorphism when rounding succeeds.
It does not prove non-isomorphism when rounding fails.
For benchmark labels, exact tools such as VF2, nauty, or bliss are used as
reference implementations.

\section{Benchmarking}

\subsection{Large Benchmark Script}

\texttt{big\_benchmark.py} compares four algorithms:
\begin{itemize}[leftmargin=*]
    \item \texttt{ours}: QP relaxation plus hyperplane rounding.
    \item \texttt{vf2}: NetworkX VF2.
    \item \texttt{nauty}: \texttt{pynauty}.
    \item \texttt{bliss}: igraph bliss.
\end{itemize}

The server benchmark configuration uses:
\begin{itemize}[leftmargin=*]
    \item 10 isomorphic pairs per graph size.
    \item 15 random pairs per graph size.
    \item graph sizes between 20 and 150 by default.
    \item random density sampled between 40\% and 85\%.
\end{itemize}

The benchmark writes resumable per-run JSON files under
\texttt{benchmark\_data/runs}.
This is necessary because the QP path is much slower than VF2, nauty, and bliss.

\subsection{Observed Runtime Pattern}

The exact tools are extremely fast on the dense random cases used so far.
The QP solver dominates runtime and memory.

A practical observation from the server run was that around $n=50$, the QP path
was already taking about 11--12 seconds per random pair.
Scaling to $n=100$ or beyond becomes expensive because the model contains
$n^2$ variables and builds $O(n^3)$ linear expression terms inside the quadratic
adjacency residual.

This is currently the main engineering bottleneck.

\section{Failure Modes}

\subsection{Diffuse Relaxed Solutions}

Graphs with symmetry or weak objective identifiability can produce diffuse
solutions.
For example, a cycle graph can allow the uniform matrix
\begin{equation}
    X_{ij} = \frac{1}{n}
\end{equation}
to satisfy the relaxed objective extremely well.
That matrix contains no sharp permutation information.

The deeper issue is not only symmetry.
The objective may have multiple equivalent or near-equivalent minimizers.
When the degree cost is weak and the adjacency term has a wide valley, rounding
can fail even if an exact permutation exists.

\subsection{Weak Cost Matrix}

Degree and neighbor-degree summaries are cheap but incomplete.
Regular graphs, strongly regular graphs, and many structured families can share
the same low-order profiles.
On those cases, $C$ may add little or no information.

A natural next phase is to add stronger node features such as Weisfeiler-Lehman
color refinement features.

\subsection{Term Scaling}

The linear degree term and quadratic adjacency term do not scale the same way.
The code now normalizes $C$, but a full normalization strategy should also
consider the typical magnitude of $\|AX-XB\|_F^2$ as $n$ changes.

\subsection{Rounding Is Heuristic}

Hyperplane rounding can miss a valid isomorphism even when one exists.
The result is then inconclusive, not a proof of non-isomorphism.

\subsection{Benchmark Labels}

Random independently sampled graph pairs are almost surely non-isomorphic for
moderate $n$, but not guaranteed.
Small-$n$ experiments must verify labels with exact tools before drawing
accuracy conclusions.

\section{What Phase 1 Establishes}

Phase 1 establishes a working and defensible pipeline:
\begin{enumerate}[leftmargin=*]
    \item the GI condition is represented as $AP=PB$;
    \item the permutation search is relaxed to the Birkhoff polytope;
    \item the implemented objective combines degree-profile compatibility with
    adjacency commutation;
    \item the relaxed solution is rounded into a candidate permutation;
    \item the final positive decision is certified exactly by $AP^*=P^*B$.
\end{enumerate}

It also establishes the limits of the current version:
\begin{enumerate}[leftmargin=*]
    \item exact solvers are much faster on the tested random dense cases;
    \item the QP construction is the runtime bottleneck;
    \item degree features are useful but weak;
    \item rounding quality is still an open research component.
\end{enumerate}

\section{Next Additions}

The clean next additions are:
\begin{itemize}[leftmargin=*]
    \item WL-based node compatibility features;
    \item stronger scaling between objective terms;
    \item Hungarian-vs-hyperplane rounding comparison on the same saved $X^*$;
    \item structured hard graph families beyond dense random graphs;
    \item a faster model construction path for the adjacency residual.
\end{itemize}

\section{Conclusion}

The current project should be presented as a research prototype, not as a
replacement for VF2, nauty, or bliss.
Its value is that it makes graph isomorphism visible as an optimization problem:
the soft matrix $X^*$, the objective components, and the rounded certificate all
expose where the method succeeds or fails.

The strongest Phase 1 claim is therefore precise:
adding the degree-profile term improves the relaxed objective's separation over
an adjacency-only baseline, while exact certification remains dependent on the
rounding stage.

\end{document}
